% Scratch research section. All labels carry the prefix tri:.
% Requires amsmath, amssymb, amsthm; no custom macros assumed.
% The face I_k is the ordinary perspective face used in the incoming reports.

\section{Mixed mass derivatives and the three-cluster polylogarithm dictionary}
\label{tri:sec:main}

The incoming two-cluster theorem evaluates
$I_{r+s}(G_a;B^r,C^s)$, and explicitly leaves $B^r,C,D$ with two
unequal displaced masses as its next target \cite{GammaColored}. We resolve that target
by a finite dictionary of weight $r+2$ and depth at most three,
uniformly in $r$. An accompanying theorem is stronger at the
differential level: the mixed derivative in arbitrarily many singly
labelled displaced masses is a divided difference of one classical
polylogarithm. The spectral sum in the Gamma application remains an
ordinary infinite sum; no arithmetic closure or minimal-depth claim
is made.

\subsection{Normalization and the inherited coefficient identity}

For $k\ge2$, write $\Sigma_k=\{t_i\ge0:\sum_i t_i=1\}$ with
measure $dt_1\cdots dt_{k-1}$. For nonempty $T\subseteq\{1,\ldots,k\}$
let $t_T=\sum_{i\in T}t_i$ and
$b_T=\sum_{i\in T}t_i b_i/t_T$. The perspective $t_TF(b_T)$ is
filled by zero when $t_T=0$. Set
\begin{align}
 \mathcal H_k(F;t)&=\sum_{\varnothing\ne T}
                (-1)^{k-|T|}t_TF(b_T),\label{tri:eq:H}\\
 I_k(F;b_1,\ldots,b_k)&=\int_{\Sigma_k}
                   \mathcal H_k(F;t)\frac{dt}{t_1\cdots t_k}.
                   \label{tri:eq:I}
\end{align}
These integrals are ordinary and absolutely convergent for $C^1$
functions on the mass interval. The perspectives are Lipschitz and
the alternating numerator vanishes on each coordinate face, so it
is bounded by a constant times $\min_i t_i$. This also supplies
locally integrable bounds after sufficiently small independent
complex powers of the $t_i$ are inserted. Constants and affine
functions are annihilated.

Let $D_F(\boldsymbol\alpha;\boldsymbol b)$ be the normalized
Dirichlet mean of $F(\sum t_i b_i)$, with total parameter
$A=\sum_i\alpha_i$. The independent-regulator identity inherited
from the two-cluster report \cite{GammaColored} implies
\begin{equation}\label{tri:eq:coefficient}
 I_k(F;\boldsymbol b)
  =[\alpha_1\cdots\alpha_k]\{A D_F(\boldsymbol\alpha;\boldsymbol b)\}.
\end{equation}
Indeed, the normalized alternating Dirichlet integral equals
$(\prod_i\alpha_i)^{-1}$ times the subset alternating sum of
$A_TD_F(\boldsymbol\alpha_T;\boldsymbol b_T)$; every proper subset
omits a variable and disappears under squarefree coefficient
extraction. When $r$ labels have the common mass $B$, their
parameters aggregate into $X$, and hence
\begin{equation}\label{tri:eq:cluster-coefficient}
 I_{r+m}(F;B^r,C_1,\ldots,C_m)
  =r![X^rY_1\cdots Y_m]\{A D_F(X,Y_1,\ldots,Y_m)\}.
\end{equation}
Here $r\ge1$, $m\ge1$, and $A=X+Y_1+\cdots+Y_m$.
All coefficients below are taken in the holomorphic continuation
of the displayed normalized mean. For the logarithms used below,
its existence also follows directly by expanding around the largest
positive spectral factor: the resulting Dirichlet moment series
has a geometric convergence margin and denominators
$(1+A)_{n-1}$.

\subsection{Arbitrarily many mixed derivatives reduce to one classical polylogarithm}

For an analytic function $f$, its divided difference \cite{DeBoor} is denoted by
$f[C_1,\ldots,C_m]$; thus one node means $f(C_1)$ and two nodes
mean $(f(C_1)-f(C_2))/(C_1-C_2)$. Repeated nodes are assigned
their unique confluent values, so
$f[C,\ldots,C]=f^{(m-1)}(C)/(m-1)!$.

\begin{theorem}[Classical-polylogarithm mixed derivatives]
\label{tri:thm:mixed}
Let $\sigma>0$, $B,C_1,\ldots,C_m\ge0$, $r,m\ge1$, and put
\begin{equation}\label{tri:eq:g}
 g_{r,\sigma,B}(C)=\frac{1}{\sigma+C}
       \operatorname{Li}_r\!\left(\frac{C-B}{\sigma+C}\right).
\end{equation}
Then
\begin{equation}\label{tri:eq:mixed}
 \boxed{\displaystyle
 \partial_{C_1}\cdots\partial_{C_m}
 I_{r+m}\bigl(U\mapsto-\log(\sigma+U);
                  B^r,C_1,\ldots,C_m\bigr)
 =(-1)^r r!\,g_{r,\sigma,B}[C_1,\ldots,C_m].}
\end{equation}
The derivative is taken in each independently labelled displaced
mass before any nodes are made equal. In particular, the bound of
one on classical polylogarithm depth is independent of both $r$ and
$m$.
\end{theorem}

\begin{proof}
Write $S_B=\sigma+B$ and $S_j=\sigma+C_j$. Differentiate the
Dirichlet mean in \eqref{tri:eq:cluster-coefficient}. Since
\[
 \partial_{C_1}\cdots\partial_{C_m}
   [-\log(S_Bt_0+\sum_jS_jt_j)]
  =(-1)^m(m-1)!\frac{t_1\cdots t_m}
                      {(S_Bt_0+\sum_jS_jt_j)^m},
\]
the coefficient of $Y_1\cdots Y_m$, first at $\Re X>0$, is
\[
 (-1)^m(m-1)!X\int_{\Sigma_{m+1}}
        \frac{t_0^{X-1}}{(S_Bt_0+\sum_jS_jt_j)^m}\,dt.
\]
This follows either by shifting each $Y_j$ by one in the Dirichlet
density or directly from $1/\Gamma(Y_j)=Y_j+O(Y_j^2)$.
At a fixed $t=t_0$, the ordinary simplex resolvent integral gives
\[
 \int_{\Sigma_m}\frac{du}
       {(S_Bt+(1-t)\sum_jS_ju_j)^m}
  =\frac{1}{(m-1)!\prod_{j=1}^m(S_Bt+(1-t)S_j)}.
\]
For completeness, this is the positive-parameter Feynman identity:
write the product of $m$ reciprocal positive numbers as a Laplace
integral over the positive orthant and separate its radial variable.
Therefore the required coefficient of the differentiated mean is
\begin{equation}\label{tri:eq:parameter-integral}
 (-1)^mX\int_0^1
  \frac{t^{X-1}(1-t)^{m-1}}
       {\prod_j(S_Bt+(1-t)S_j)}\,dt.
\end{equation}
At fixed $t$, the divided difference of
$C\mapsto(S_Bt+(1-t)(\sigma+C))^{-1}$ is
$(-1)^{m-1}(1-t)^{m-1}$ divided by the product in
\eqref{tri:eq:parameter-integral}. Thus that expression is the
negative divided difference of
\[
 R_X(C)=X\int_0^1
             \frac{t^{X-1}}{S_Bt+(1-t)(\sigma+C)}\,dt.
\]
Writing $z=(C-B)/(\sigma+C)$ yields, initially for $|z|<1$,
\begin{equation}\label{tri:eq:RX}
 R_X(C)=\frac1{\sigma+C}
       \left(1+\sum_{n=1}^{\infty}\frac{Xz^n}{X+n}\right),
 \qquad
 [X^r]R_X(C)=\frac{(-1)^{r-1}}{\sigma+C}\operatorname{Li}_r(z).
\end{equation}
The series is normally convergent near $X=0$ on each compact
$z$-disk. Its integral representation continues it through every
real $z<1$ without crossing a singularity; equivalently, subtract
the value at $t=0$ before continuing the $X$ integral. Now take
$r![X^r]$ and use \eqref{tri:eq:cluster-coefficient}. All mass
derivatives and confluent limits are justified by local holomorphy
away from nonpositive spectral factors. This proves the theorem.
\end{proof}

For $m=1$, integration from $C=B$ recovers the inherited
one-displaced formula
\begin{equation}\label{tri:eq:one-displaced}
 I_{r+1}(-\log(\sigma+U);B^r,C)
   =(-1)^r r!\operatorname{Li}_{1,r}
                         \!\left(\frac{C-B}{\sigma+C}\right),
\end{equation}
because $d\operatorname{Li}_{1,r}(z)/dz=
\operatorname{Li}_r(z)/(1-z)$ and the face vanishes when all masses
coincide. This recovery checks the factorial and sign conventions.

\subsection{A uniform finite dictionary for the three-cluster primitive}

It is useful to specify colored multiple polylogarithms by cumulative
colors, which avoids apparent quotients by zero. For positive
indices $s_1,\ldots,s_d$, define
\begin{equation}\label{tri:eq:cumulative}
 \mathcal C_{s_1,\ldots,s_d}(c_1,\ldots,c_d)
 =\sum_{n_1>\cdots>n_d\ge1}
 \frac{c_1^{n_1-n_2}c_2^{n_2-n_3}\cdots
            c_{d-1}^{n_{d-1}-n_d}c_d^{n_d}}
      {n_1^{s_1}\cdots n_d^{s_d}}.
\end{equation}
For $|c_j|<1$ this series converges absolutely. In the conventional
colored notation it is
$\operatorname{Li}_{s_1,\ldots,s_d}
(c_1,c_2/c_1,\ldots,c_d/c_{d-1})$, interpreted continuously
when a color is zero. The one-variable convention used by the
incoming reports is
$\operatorname{Li}_{s_1,\ldots,s_d}(z)=
\mathcal C_{s_1,\ldots,s_d}(z,\ldots,z)$.
For real cumulative colors below one, the continuation is the
iterated-integral branch continued from zero, with letters
$0,1/c_1,\ldots,1/c_d$ and the path $[0,1]$. No letter crosses
that path on the stated real domain.

Put $H(z)=\operatorname{Li}_1(z)=-\log(1-z)$ and
\begin{align}
 \mathcal S_r(z,w)
   ={}&\mathcal C_{1,1,r}(z,w,w)
      +\mathcal C_{1,r,1}(z,z,w)
      +\mathcal C_{1,r+1}(z,w),\label{tri:eq:S}\\
 \mathcal K_r(z,w)
   ={}&H(w)\operatorname{Li}_{1,r}(z)
       +H(z)\operatorname{Li}_{1,r}(w)
       -\mathcal S_r(z,w)-\mathcal S_r(w,z),\label{tri:eq:K}\\
 \mathcal P_r(z,w)
   ={}&(r+1)\{\operatorname{Li}_{1,r+1}(z)
                    +\operatorname{Li}_{1,r+1}(w)\}
       +\mathcal K_r(z,w).\label{tri:eq:P}
\end{align}
Every term has weight $r+2$; the products in \eqref{tri:eq:K}
can be shuffled as iterated integrals into depth at most three,
retaining the same safe letters for real negative colors.
Thus the formula gives
a finite dictionary with an upper bound of three on depth for
every $r\ge1$.

\begin{theorem}[Three-cluster logarithmic face]
\label{tri:thm:three}
For $\sigma>0$, $B,C,D\ge0$, $r\ge1$, let
$z=(C-B)/(\sigma+C)$ and $w=(D-B)/(\sigma+D)$. Then
\begin{equation}\label{tri:eq:three}
 \boxed{\displaystyle
 I_{r+2}(-\log(\sigma+U);B^r,C,D)
       =(-1)^{r+1}r!\mathcal P_r(z,w).}
\end{equation}
In particular, if $B\le C,D$, all cumulative colors in the finite
dictionary lie in $[0,1)$, regardless of the sizes of the mass
displacements. The result therefore requires no small-displacement
hypothesis.
\end{theorem}

\begin{proof}
The key harmonic identity, valid for $|z|,|w|<1$, is
\begin{equation}\label{tri:eq:K-series}
 \mathcal K_r(z,w)
  =\sum_{m,n\ge1}\frac{z^mw^n}{mn}
    \{H_{m-1}^{(r)}+H_{n-1}^{(r)}-H_{m+n-1}^{(r)}\}.
\end{equation}
To check its finite reduction, the first two harmonic terms give
the two products in \eqref{tri:eq:K}. For the last term use
$1/(mn)=(1/m+1/n)/(m+n)$ and put $N=m+n$. The harmonic index
$j<N$ can lie below $m$, above $m$, or equal $m$; these three
cases give respectively
$\mathcal C_{1,1,r}(z,w,w)$,
$\mathcal C_{1,r,1}(z,z,w)$, and
$\mathcal C_{1,r+1}(z,w)$. Exchanging $z,w$ supplies the other
half. All rearrangements are absolute.

Write $f_r(z)=(1-z)\operatorname{Li}_r(z)$. Differentiating
\eqref{tri:eq:K-series} gives
\begin{equation}\label{tri:eq:K-mixed}
 \partial_z\partial_w\mathcal K_r(z,w)
   =-\frac{f_r[z,w]}{(1-z)(1-w)}.
\end{equation}
Indeed,
$\sum_{n\ge1}H_{n-1}^{(r)}z^{n-1}=
\operatorname{Li}_r(z)/(1-z)$, while the diagonal sum involving
$H_{m+n-1}^{(r)}$ equals
\[
 \frac{\operatorname{Li}_r(z)/(1-z)
                    -\operatorname{Li}_r(w)/(1-w)}{z-w}.
\]
Combining these expressions proves \eqref{tri:eq:K-mixed},
including its removable diagonal.

Now $z'=(1-z)/(\sigma+C)$,
$w'=(1-w)/(\sigma+D)$,
\[
 z-w=\frac{(\sigma+B)(C-D)}{(\sigma+C)(\sigma+D)},
 \qquad f_r(z)=(\sigma+B)g_{r,\sigma,B}(C).
\]
Consequently the mixed $C,D$ derivative of the right side of
\eqref{tri:eq:three} agrees exactly with
Theorem~\ref{tri:thm:mixed} for $m=2$. On $C=B$ or $D=B$,
$\mathcal K_r$ vanishes and the remaining term is precisely
the inherited one-displaced formula with $r+1$ base labels.
Equality follows from two integrations of the common mixed
derivative, starting on those two coordinate edges. The proof
initially applies in a neighborhood of $C=D=B$ and throughout
$B\le C,D$ by the absolutely convergent color series. Analytic
continuation along the positive mass domain gives every stated
$B,C,D\ge0$.
\end{proof}

The formula is symmetric under $C\leftrightarrow D$. Its two
principal collision checks are
\begin{align}
 \mathcal P_r(z,0)&=(r+1)\operatorname{Li}_{1,r+1}(z),
                      \label{tri:eq:axis}\\
 \mathcal P_r(z,z)&=2r\operatorname{Li}_{1,r+1}(z)
       +2\sum_{j=1}^{r}\operatorname{Li}_{1,j,r+1-j}(z).
                      \label{tri:eq:diagonal}
\end{align}
The second identity recovers twice the predecessor's $P_{r,2}$.
It also follows directly from the shuffle identity
\[
 H(z)\operatorname{Li}_{1,r}(z)
  =\operatorname{Li}_{1,1,r}(z)
      +\operatorname{Li}_{1,r,1}(z)
      +\sum_{j=1}^{r}\operatorname{Li}_{1,j,r+1-j}(z).
\]
Thus the three-cluster formula has the correct labelled
multiplicity even when the two displaced masses become equal.

\subsection{Gamma faces, confluent derivatives, and polygamma values}

For $a>0$, $U\ge0$, put
\[
 G_a(U)=\log\Gamma(a+i\sqrt U)+\log\Gamma(a-i\sqrt U)
                                     -\log(2\pi).
\]
The analytic logarithms are continued from positive real Gamma
arguments. The ordinary Gamma product \cite[\S5.8]{DLMF} gives
\[
 G_a(U)=2\log\Gamma(a)-\log(2\pi)
       -\sum_{n=0}^{\infty}
           \log\left(1+\frac{U}{(n+a)^2}\right).
\]
Its second mass derivative is normally convergent on compact
nonnegative mass sets, so every ordinary face may be applied
termwise after subtracting an affine function from each term.

\begin{corollary}[The three-cluster Gamma dictionary]
\label{tri:cor:Gamma}
For $a>0$, $B,C,D\ge0$ and $r\ge1$, set
\[
 z_n=\frac{C-B}{(n+a)^2+C},\qquad
 w_n=\frac{D-B}{(n+a)^2+D}.
\]
Then
\begin{equation}\label{tri:eq:Gamma}
 \boxed{\displaystyle
 I_{r+2}(G_a;B^r,C,D)
   =(-1)^{r+1}r!\sum_{n=0}^{\infty}\mathcal P_r(z_n,w_n).}
\end{equation}
The sum is ordinary and absolutely convergent, locally normally
in the stated parameter domain. Each spectral term has a finite
colored multiple-polylogarithm representation of depth at most
three.
\end{corollary}

\begin{proof}
Apply Theorem~\ref{tri:thm:three} termwise to the Gamma product.
For a direct tail bound, \eqref{tri:eq:K-series} gives
\[
 \mathcal P_r(z,w)=\frac{r+1}{2}(z^2+w^2)-zw
                     +O((|z|+|w|)^3).
\]
Both colors are $O(n^{-2})$, so the spectral summands are
$O(n^{-4})$. Compact local holomorphic extensions give the same
normal convergence for every fixed parameter derivative.
\end{proof}

Define the normally convergent classical-polylogarithm sum
\begin{equation}\label{tri:eq:spectralg}
 \mathfrak g_{r,a,B}(C)=\sum_{n=0}^{\infty}
      \frac{\operatorname{Li}_r((C-B)/((n+a)^2+C))}
           {(n+a)^2+C}.
\end{equation}
Theorem~\ref{tri:thm:mixed} gives, for every $m\ge1$,
\begin{equation}\label{tri:eq:Gamma-mixed}
 \boxed{\displaystyle
 \partial_{C_1}\cdots\partial_{C_m}
 I_{r+m}(G_a;B^r,C_1,\ldots,C_m)
   =(-1)^r r!\mathfrak g_{r,a,B}[C_1,\ldots,C_m].}
\end{equation}
For example, the first case beyond the two-cluster primitive is
\begin{equation}\label{tri:eq:dilog-example}
 \partial_C\partial_D I_4(G_a;B,B,C,D)
   =\frac{2}{C-D}\sum_{n=0}^{\infty}
      \left\{\frac{\operatorname{Li}_2(z_n)}{(n+a)^2+C}
             -\frac{\operatorname{Li}_2(w_n)}{(n+a)^2+D}\right\}.
\end{equation}
At $C=D$ it is the corresponding derivative of
$\mathfrak g_{2,a,B}$; no singular mass denominator remains.

More generally, making $m$ independent displaced masses equal
after differentiation gives
\begin{equation}\label{tri:eq:confluent}
 \left.\partial_{C_1}\cdots\partial_{C_m}
 I_{r+m}(G_a;B^r,C_1,\ldots,C_m)\right|_{C_j=C}
 =\frac{(-1)^r r!}{(m-1)!}\,
                    \partial_C^{m-1}\mathfrak g_{r,a,B}(C).
\end{equation}
Every finite derivative here involves only classical polylogarithms
$\operatorname{Li}_{r-j}$ and rational mass factors. Negative
polylogarithm orders are rational functions. The distinction
between these labelled mixed derivatives and total derivatives
along a diagonal is essential.

At the complete mass collision $C=B$, the formula has a finite
harmonic-polygamma specialization. Let
$H_n^\star(\{1\}^r)=\sum_{n\ge j_1\ge\cdots\ge j_r\ge1}
(j_1\cdots j_r)^{-1}$. For $m\ge2$,
\begin{equation}\label{tri:eq:complete-collision}
 \left.\partial_{C_1}\cdots\partial_{C_m}
 I_{r+m}(G_a;B^r,C_1,\ldots,C_m)\right|_{C_j=B}
 =(-1)^{r+m}r!H_{m-1}^\star(\{1\}^r)
             \sum_{n=0}^{\infty}((n+a)^2+B)^{-m}.
\end{equation}
Indeed the coefficient of $u^N$ in
$(1+u)^{-1}\operatorname{Li}_r(u/(1+u))$ is
$(-1)^{N+1}H_N^\star(\{1\}^r)$; this follows from the elementary
binomial identity
$\sum_{j=1}^N(-1)^{j-1}\binom Nj/j^r=H_N^\star(\{1\}^r)$.
At $B=0$ the final sum is $\zeta(2m,a)$; at $B>0$ it is a
finite derivative of the paired digamma resolvent
\[
 \sum_{n=0}^{\infty}((n+a)^2+B)^{-1}
  =\frac{\psi(a+i\sqrt B)-\psi(a-i\sqrt B)}{2i\sqrt B}.
\]
This is a consistency check with the predecessor's all-degree
collision formula, not a separate novelty claim for those
coefficients.

\subsection{An independent ordinary integral for verification}

The finite dictionary also admits a single scalar logarithmic
integral that is convenient for numerical verification. Put
$q=(B-C)/(\sigma+B)$, $p=(B-D)/(\sigma+B)$,
$L(t)=\log(1-t)$, $a(t)=-\log(1-qt)$, and
$b(t)=-\log(1-pt)$. Then
\begin{align}\label{tri:eq:log-integral}
 I_{r+2}(-\log(\sigma+U);B^r,C,D)
  =\int_0^1\bigg[&L^r\{a'(L+b)+b'(L+a)\}\notag\\
       &\quad +(r+1)\frac{L^r(a+b)}t
       +r\frac{L^{r-1}ab}t\bigg]dt.
\end{align}
All apparent singularities at zero are removable or integrable;
at one only powers of $\log(1-t)$ occur. One derivation uses
independent Gamma variables and Frullani's formula to write
\[
 A D_{-\log}(X,Y,Z)
   =\mathcal L(X,Y,Z)
     +\int_0^1(1-t)^{X+Y+Z}
       \frac{d}{dt}\frac{(1-qt)^{-Y}(1-pt)^{-Z}-1}{t}\,dt.
\]
Here $\mathcal L$ is linear in $X,Y,Z$.
Extract $r![X^rYZ]$ and integrate the double-pole term using
the primitive $(1-t)/t$, which vanishes at $t=1$.
This produces \eqref{tri:eq:log-integral} without endpoint
finite parts.

For an additional check wholly in the local mass series, let
$h_j(N)=H_{N-1}^{\star}(\{1\}^j)$. When $|q|,|p|<1$,
the independent Dirichlet moment expansion gives
\begin{align}\label{tri:eq:collision-check}
 I_{r+2}(-\log(\sigma+U);B^r,C,D)
 =(-1)^r r!\bigg[&\sum_{m,n\ge1}\frac{q^mp^n}{mn}h_r(m+n)\notag\\
 &-(r+1)\sum_{m\ge1}\frac{q^m+p^m}{m}h_{r+1}(m)\bigg].
\end{align}
The two sums arise respectively from terms using both displaced
Dirichlet parameters and from terms using just one. This provides
an independently organized harmonic computation of the same face.

The supplied verifier records 231 assertions: 192 exact rational
binomial--harmonic checks, fourteen comparisons from seven scalar
parameter cases, 24 independently computed mixed-derivative checks,
and one Gamma spectral check. The scalar cases include
$r=1,2,3,6$, a displaced mass below the base, an equal-outlier
collision, and a zero displacement. At $\sigma=1$, $B=0$,
$C=1/5$, $D=2/5$, $r=2$, for example, the common scalar value is
\[
 -0.307903203351030388542590020916578537605405701313403\ldots.
\]
The finite colored dictionary, the ordinary scalar integral, and
the local Dirichlet series agree at 75-digit working precision
with scaled residuals below $1.5\cdot10^{-75}$. The independent
mixed-derivative checks use the predecessor's monomial-face
formula and include both complete and partial confluence; their
largest scaled residual is below $6\cdot10^{-72}$.
The Gamma example is
\begin{equation}\label{tri:eq:Gamma-value}
 I_4(G_1;0,0,1/5,2/5)
  =-0.3411810552611179668555780745233632208321622014483386\ldots.
\end{equation}
A finite colored-MPL spectral head plus an exact-form Hurwitz
tail agrees with the independent complete Hurwitz collision
series to a scaled residual below $1.4\cdot10^{-76}$.
These are floating-point diagnostics, not interval certificates
or substitutes for the analytic proofs.

